\section{Estructura del curso, tareas y proyectos}

\subsection{Estratificación del contenido del curso}

\begin{figure}[H]
\centering
\includegraphics[page=11,width=0.9\textwidth]{slides.pdf}
\caption{Este curso divide el contenido en cuatro líneas: capacidades del modelo, Agent framework, aplicaciones y seguridad y ética.}
\end{figure}

La estructura del curso se divide en cuatro niveles:
\begin{enumerate}
    \item \textbf{Model core capabilities}: razonamiento, planeación, comprensión multimodal.
    \item \textbf{LLM agent frameworks}: workflow, tool use, RAG, multi-Agent.
    \item \textbf{Applications}: desarrollo de software, automatización de flujos de trabajo, aplicaciones multimodales y empresariales.
    \item \textbf{Safety \& ethics}: riesgos de despliegue, interacción humano-máquina, problemas de privacidad y control.
\end{enumerate}

\begin{knowledgebox}{Intención de diseño de las direcciones del proyecto}
El proyecto del curso no busca que los estudiantes repitan la construcción de un chatbot, sino que fomenta desarrollarlo en torno a cinco tipos de temas: aplicaciones, puntos de referencia, capacidades fundamentales, seguridad, multi-Agent o sistemas descentralizados. Este diseño en sí mismo indica a los estudiantes que la investigación en Agentes no es una dirección única, sino un espacio de investigación conformado en conjunto por capacidades, sistemas y gobernanza.
\end{knowledgebox}

\subsection{Evaluación y organización del proyecto}

La introducción también presenta con claridad la manera en que se estructura la entrega del curso:
\begin{itemize}
    \item tareas de lectura semanales, que sirven para construir un contexto común entre las distintas clases;
    \item un hands-on lab, para poner realmente en marcha los frameworks y las herramientas;
    \item un semester-long project, que fomenta formar equipos de cinco integrantes y desarrollarlo en torno a un problema completo.
\end{itemize}

En cuanto a los requisitos de calidad, el curso hace corresponder de 1 a 4 unit con distintos niveles de profundidad de implementación, lo cual muestra que da más importancia a «convertir el sistema de Agentes en una entidad de ingeniería que se pueda explicar, mostrar y verificar», y no solo a leer artículos o escribir reseñas conceptuales.

\subsection{Resumen de la sección}

La manera en que se organiza este curso es muy congruente con su propio tema: el curso mismo es un campo de entrenamiento de sistemas de Agentes que va de la lectura y la experimentación hasta la implementación del proyecto.

